\documentclass[../../main.tex]{subfiles}% 注意这里的文件路径不能用 ./main.tex ,否则用latexmk编译子文件会报错
\graphicspath{{\subfix{./image/}}{\subfix{../../image/}}} % 指定图片引用路径,后续可以直接使用图片文件名
% 如果图片存放在上级目录,则使用 ../../image/ 作为备用路径

% 例如:
% \begin{figure}[H]
% \centering
% \includegraphics[width=0.3\textwidth]{图.png}
% \caption{}
% \label{figure:图}
% \end{figure}
% 注意:上述\label{}一定要放在\caption{}之后,否则引用图片序号会只会显示??.

\begin{document}

% SOURCE p=28; printed=13; object=ch00-sec06-title; region=top
\section{度量空间}\label{section:ch00-metric-spaces}

% SOURCE p=28; printed=13; object=ch00-metric-definition; region=upper

\begin{definition}[度量与度量空间]\label{definition:ch00-metric-definition}
集合 $X$ 上的\textbf{度量}是满足下列条件的函数 $\rho:X\times X\to[0,\infty)$：
\begin{itemize}
\item $\rho(x,y)=0$ 当且仅当 $x=y$；
\item 对所有 $x,y\in X$，$\rho(x,y)=\rho(y,x)$；
\item 对所有 $x,y,z\in X$，$\rho(x,z)\leqslant\rho(x,y)+\rho(y,z)$。
\end{itemize}
（直观地，$\rho(x,y)$ 应理解为从 $x$ 到 $y$ 的距离。）配备了一个度量的集合称为\textbf{度量空间}。
\end{definition}
下面给出几个例子。

% SOURCE p=28; printed=13; object=ch00-metric-examples; region=middle

\begin{examples}\label{examples:ch00-metric-examples}
\begin{enumerate}[label=\roman*.,ref=(\roman*)]
\item Euclid 距离 $\rho(x,y)=|x-y|$ 是 $\mathbb{R}^n$ 上的度量。
\item $\rho_1(f,g)=\int_0^1|f(x)-g(x)|\,\mathrm{d}x$ 和 $\rho_\infty(f,g)=\sup_{0\leqslant x\leqslant1}|f(x)-g(x)|$ 是 $[0,1]$ 上连续函数空间上的度量。
\item 若 $\rho$ 为 $X$ 上的度量且 $A\subseteq X$，则 $\rho|_{A\times A}$ 是 $A$ 上的度量。
\item 
\begin{definition*}[乘积度量]
若 $(X_1,\rho_1)$ 和 $(X_2,\rho_2)$ 是度量空间，则 $X_1\times X_2$ 上的\textbf{乘积度量} $\rho$ 定义为
\begin{align*}
\rho\bigl((x_1,x_2),(y_1,y_2)\bigr)=\max\bigl(\rho_1(x_1,y_1),\rho_2(x_2,y_2)\bigr).
\end{align*}
\end{definition*}

有时也在 $X_1\times X_2$ 上使用其他度量，例如
\begin{align*}
\rho_1(x_1,y_1)+\rho_2(x_2,y_2)\quad\text{或}\quad
\bigl[\rho_1(x_1,y_1)^2+\rho_2(x_2,y_2)^2\bigr]^{1/2}.
\end{align*}
不过，在本节末将定义的意义下，这些度量与乘积度量是等价的。
\end{enumerate}
\end{examples}


% SOURCE p=28; printed=13; object=ch00-balls-open-closed; region=lower

\begin{definition}[开球、开集与闭集]\label{definition:ch00-balls-open-closed}
设 $(X,\rho)$ 为度量空间。若 $x\in X$ 且 $r>0$，以 $x$ 为中心、半径为 $r$ 的（开）\textbf{球}为
\begin{align*}
B(r,x)=\{y\in X:\rho(x,y)<r\}.
\end{align*}
集合 $E\subseteq X$ 为\textbf{开集}，是指对每个 $x\in E$，存在 $r>0$ 使 $B(r,x)\subseteq E$；若其补集为开集，则称 $E$ 为\textbf{闭集}。
\end{definition}

\begin{example}\label{example:ch00-balls-open-example}
例如，每个球 $B(r,x)$ 都是开集，因为若 $y\in B(r,x)$ 且 $\rho(x,y)=s$，则 $B(r-s,y)\subseteq B(r,x)$。此外，$X$ 和 $\varnothing$ 既开又闭。
\end{example}

\begin{proposition}\label{proposition:ch00-open-closed-operations}
显然，任意开集族的并为开集，从而任意闭集族的交为闭集。另外，有限个开集（分别地，闭集）的交（分别地，并）为开集（分别地，闭集）。
\end{proposition}

\begin{proof}
事实上，若 $U_1,\ldots,U_n$ 为开集且 $x\in\bigcap_1^nU_j$，则对每个 $j$，存在 $r_j>0$ 使 $B(r_j,x)\subseteq U_j$；于是 $B(r,x)\subseteq\bigcap_1^nU_j$，其中 $r=\min(r_1,\ldots,r_n)$，所以 $\bigcap_1^nU_j$ 为开集。
\end{proof}


% SOURCE p=28; printed=13; object=ch00-interior-closure; region=bottom

\begin{definition}[内部、闭包与稠密性]\label{definition:ch00-interior-closure}
若 $E\subseteq X$，则所有满足 $U\subseteq E$ 的开集 $U$ 的并，是包含于 $E$ 的最大开集，称为 $E$ 的\textbf{内部}，记为 $E^\circ$。类似地，所有满足 $F\supseteq E$ 的闭集 $F$ 的交，是包含 $E$ 的最小闭集，称为 $E$ 的\textbf{闭包}，记为 $\overline{E}$。若 $\overline{E}=X$，则称 $E$ 在 $X$ 中\textbf{稠密}；
% SOURCE p=29; printed=14; object=ch00-density-separability-convergence; region=top
若 $\overline{E}$ 的内部为空，则称 $E$ \textbf{无处稠密}。若 $X$ 有可数稠密子集，则称 $X$ \textbf{可分}。
\end{definition}

\begin{example}\label{example:ch00-rational-dense-example}
（例如，$\mathbb{Q}^n$ 是 $\mathbb{R}^n$ 的可数稠密子集。）
\end{example}

\begin{definition}[序列的收敛]\label{definition:ch00-metric-convergence}
若 $\lim_{n\to\infty}\rho(x_n,x)=0$，则称 $X$ 中的序列 $\{x_n\}$ \textbf{收敛}于 $x\in X$，记为 $x_n\to x$ 或 $\lim x_n=x$。
\end{definition}


% SOURCE p=29; printed=14; object=ch00-closure-characterization; number=0.22; region=upper
\begin{proposition}\label{proposition:ch00-closure-characterization}
若 $X$ 为度量空间，$E\subseteq X$，且 $x\in X$，则下列条件等价：
\begin{enumerate}[label=\alph*.,ref=(\alph*)]
\item\label{item:ch00-closure-member} $x\in\overline{E}$。
\item\label{item:ch00-closure-balls} 对所有 $r>0$，$B(r,x)\cap E\ne\varnothing$。
\item\label{item:ch00-closure-sequence} 存在 $E$ 中的序列 $\{x_n\}$ 收敛于 $x$。
\end{enumerate}
\end{proposition}

% SOURCE p=29; printed=14; object=ch00-proof-closure-characterization; region=upper-middle
\begin{proof}
若 $B(r,x)\cap E=\varnothing$，则 $B(r,x)^{\mathrm{c}}$ 是包含 $E$ 但不包含 $x$ 的闭集，所以 $x\notin\overline{E}$。反之，若 $x\notin\overline{E}$，由于 $(\overline{E})^{\mathrm{c}}$ 为开集，存在 $r>0$ 使 $B(r,x)\subseteq(\overline{E})^{\mathrm{c}}\subseteq E^{\mathrm{c}}$。因此\subcref{proposition:ch00-closure-characterization}{item:ch00-closure-member}与\subcref{proposition:ch00-closure-characterization}{item:ch00-closure-balls}等价。若\subcref{proposition:ch00-closure-characterization}{item:ch00-closure-balls}成立，则对每个 $n\in\mathbb{N}$，存在 $x_n\in B(n^{-1},x)\cap E$，所以 $x_n\to x$。另一方面，若 $B(r,x)\cap E=\varnothing$，则对所有 $y\in E$ 都有 $\rho(y,x)\geqslant r$，所以 $E$ 中没有序列能收敛于 $x$。因此\subcref{proposition:ch00-closure-characterization}{item:ch00-closure-balls}与\subcref{proposition:ch00-closure-characterization}{item:ch00-closure-sequence}等价。
\end{proof}

% SOURCE p=29; printed=14; object=ch00-continuity-definition; region=middle

\begin{definition}[连续与一致连续]\label{definition:ch00-continuity-definition}
若 $(X_1,\rho_1)$ 和 $(X_2,\rho_2)$ 是度量空间，映射 $f:X_1\to X_2$ 在 $x\in X$ 处\textbf{连续}，是指对每个 $\varepsilon>0$，存在 $\delta>0$，使得只要 $\rho_1(x,y)<\delta$，就有 $\rho_2(f(y),f(x))<\varepsilon$；换句话说，$f^{-1}(B(\varepsilon,f(x)))\supseteq B(\delta,x)$。若 $f$ 在每个 $x\in X_1$ 处连续，则称映射 $f$ \textbf{连续}；若还能在连续性的定义中选取与 $x$ 无关的 $\delta$，则称 $f$ \textbf{一致连续}。
\end{definition}


% SOURCE p=29; printed=14; object=ch00-continuity-open-preimages; number=0.23; region=middle-lower
\begin{proposition}\label{proposition:ch00-continuity-open-preimages}
$f:X_1\to X_2$ 连续，当且仅当对每个开集 $U\subseteq X_2$，$f^{-1}(U)$ 在 $X_1$ 中为开集。
\end{proposition}

% SOURCE p=29; printed=14; object=ch00-proof-continuity-open-preimages; region=lower
\begin{proof}
若后一条件成立，则对每个 $x\in X_1$ 和 $\varepsilon>0$，集合 $f^{-1}(B(\varepsilon,f(x)))$ 是包含 $x$ 的开集，因而包含某个以 $x$ 为中心的球；这意味着 $f$ 在 $x$ 处连续。反之，设 $f$ 连续且 $U$ 在 $X_2$ 中为开集。对每个 $y\in U$，存在 $\varepsilon_y>0$ 使 $B(\varepsilon_y,y)\subseteq U$；对每个 $x\in f^{-1}(\{y\})$，存在 $\delta_x>0$ 使 $B(\delta_x,x)\subseteq f^{-1}(B(\varepsilon_y,y))\subseteq f^{-1}(U)$。因此 $f^{-1}(U)=\bigcup_{x\in f^{-1}(U)}B(\delta_x,x)$ 为开集。
\end{proof}

% SOURCE p=29; printed=14; object=ch00-completeness-definition; region=lower

\begin{definition}[Cauchy 序列与完备性]\label{definition:ch00-completeness-definition}
度量空间 $(X,\rho)$ 中的序列 $\{x_n\}$ 称为 \textbf{Cauchy 序列}，是指当 $n,m\to\infty$ 时，$\rho(x_n,x_m)\to0$。$X$ 的子集 $E$ 称为\textbf{完备的}，是指 $E$ 中每个 Cauchy 序列都收敛，且极限属于 $E$。
\end{definition}

\begin{example}\label{example:ch00-complete-space-example}
例如，$\mathbb{R}^n$（配备 Euclid 度量）是完备的，而 $\mathbb{Q}^n$ 不是。
\end{example}


% SOURCE p=29; printed=14; object=ch00-closed-complete; number=0.24; region=bottom
\begin{proposition}\label{proposition:ch00-closed-complete}
完备度量空间的闭子集是完备的，任意度量空间的完备子集是闭的。
\end{proposition}

% SOURCE p=29; printed=14; object=ch00-proof-closed-complete; region=bottom
\begin{proof}
若 $X$ 完备，$E\subseteq X$ 为闭集，且 $\{x_n\}$ 是 $E$ 中的 Cauchy 序列，则 $\{x_n\}$ 在 $X$ 中有极限 $x$。由\cref{proposition:ch00-closure-characterization}，$x\in\overline{E}=E$。若 $E\subseteq X$ 完备且 $x\in\overline{E}$，由\cref{proposition:ch00-closure-characterization}，存在 $E$ 中的序列 $\{x_n\}$ 收敛于 $x$。$\{x_n\}$ 为 Cauchy 序列，所以其极限属于 $E$；因此 $E=\overline{E}$。
\end{proof}

% SOURCE p=30; printed=15; object=ch00-distance-diameter; region=top

\begin{definition}[点与集合、集合与集合之间的距离]\label{definition:ch00-distance-diameter}
在度量空间 $(X,\rho)$ 中，可以定义点到集合的距离，以及两个集合之间的距离。具体地，若 $x\in X$ 且 $E,F\subseteq X$，则
\begin{align*}
\rho(x,E)&=\inf\{\rho(x,y):y\in E\},\\
\rho(E,F)&=\inf\{\rho(x,y):x\in E,\ y\in F\}=\inf\{\rho(x,F):x\in E\}.
\end{align*}
\end{definition}


\begin{remark}
注意，由\cref{proposition:ch00-closure-characterization}，$\rho(x,E)=0$ 当且仅当 $x\in\overline{E}$。
\end{remark}

\begin{definition}[直径与有界集]\label{definition:ch00-diameter-boundedness}
还定义 $E\subseteq X$ 的\textbf{直径}为
\begin{align*}
\operatorname{diam}E=\sup\{\rho(x,y):x,y\in E\}.
\end{align*}
若 $\operatorname{diam}E<\infty$，则称 $E$ \textbf{有界}。
\end{definition}


% SOURCE p=30; printed=15; object=ch00-total-boundedness; region=upper-middle

\begin{definition}[覆盖与全有界集]\label{definition:ch00-total-boundedness}
若 $E\subseteq X$，且 $\{V_\alpha\}_{\alpha\in A}$ 是满足 $E\subseteq\bigcup_{\alpha\in A}V_\alpha$ 的集合族，则称 $\{V_\alpha\}_{\alpha\in A}$ 为 $E$ 的一个\textbf{覆盖}，并称 $E$ 被这些 $V_\alpha$ \textbf{覆盖}。若对每个 $\varepsilon>0$，$E$ 都可被有限个半径为 $\varepsilon$ 的球覆盖，则称 $E$ \textbf{全有界}。
\end{definition}

\begin{proposition}\label{proposition:ch00-total-implies-bounded}
每个全有界集都有界。
\end{proposition}

\begin{proof}
若 $x,y\in\bigcup_1^nB(\varepsilon,z_j)$，设 $x\in B(\varepsilon,z_1)$ 且 $y\in B(\varepsilon,z_2)$，则
\begin{align*}
\rho(x,y)&\leqslant\rho(x,z_1)+\rho(z_1,z_2)+\rho(z_2,y)\\
&\leqslant2\varepsilon+\max\{\rho(z_j,z_k):1\leqslant j,k\leqslant n\}.
\end{align*}
\end{proof}


\begin{remark}
（逆命题一般不成立。）若 $E$ 全有界，则 $\overline{E}$ 也全有界，因为容易看出，若 $E\subseteq\bigcup_1^nB(\varepsilon,z_j)$，则 $\overline{E}\subseteq\bigcup_1^nB(2\varepsilon,z_j)$。
\end{remark}


% SOURCE p=30; printed=15; object=ch00-compactness-characterization; number=0.25; region=middle
\begin{theorem}\label{theorem:ch00-compactness-characterization}
若 $E$ 是度量空间 $(X,\rho)$ 的子集，则下列条件等价：
\begin{enumerate}[label=\alph*.,ref=(\alph*)]
\item\label{item:ch00-compact-complete-total} $E$ 完备且全有界。
\item\label{item:ch00-compact-bolzano-weierstrass} （\textbf{Bolzano–Weierstrass 性质}）$E$ 中的每个序列都有一个子序列收敛于 $E$ 中的某一点。
\item\label{item:ch00-compact-heine-borel} （\textbf{Heine–Borel 性质}）若 $\{V_\alpha\}_{\alpha\in A}$ 为 $E$ 的开覆盖，则存在有限集 $F\subseteq A$，使 $\{V_\alpha\}_{\alpha\in F}$ 覆盖 $E$。
\end{enumerate}
\end{theorem}

% SOURCE p=30; printed=15; object=ch00-proof-compactness-plan; region=lower-middle
\begin{proof}
我们将证明\subcref{theorem:ch00-compactness-characterization}{item:ch00-compact-complete-total}与\subcref{theorem:ch00-compactness-characterization}{item:ch00-compact-bolzano-weierstrass}等价，二者共同推出\subcref{theorem:ch00-compactness-characterization}{item:ch00-compact-heine-borel}，最后证明\subcref{theorem:ch00-compactness-characterization}{item:ch00-compact-heine-borel}推出\subcref{theorem:ch00-compactness-characterization}{item:ch00-compact-bolzano-weierstrass}。

% SOURCE p=30; printed=15; object=ch00-proof-compactness-a-b; region=lower
\subcref{theorem:ch00-compactness-characterization}{item:ch00-compact-complete-total}推出\subcref{theorem:ch00-compactness-characterization}{item:ch00-compact-bolzano-weierstrass}：设前者成立，且 $\{x_n\}$ 为 $E$ 中的序列。$E$ 可被有限个半径为 $2^{-1}$ 的球覆盖，其中至少有一个球包含无穷多个 $x_n$；设当 $n\in N_1$ 时 $x_n\in B_1$。$E\cap B_1$ 可被有限个半径为 $2^{-2}$ 的球覆盖，其中至少有一个球包含无穷多个指标属于 $N_1$ 的 $x_n$；设当 $n\in N_2$ 时 $x_n\in B_2$。归纳进行，得到半径为 $2^{-j}$ 的球序列 $B_j$，以及 $\mathbb{N}$ 的递减子集序列 $N_j$，使得当 $n\in N_j$ 时 $x_n\in B_j$。选取 $n_1\in N_1,n_2\in N_2,\ldots$，使 $n_1<n_2<\cdots$。则 $\{x_{n_j}\}$ 为 Cauchy 序列，因为当 $k>j$ 时 $\rho(x_{n_j},x_{n_k})<2^{1-j}$；由 $E$ 的完备性，它在 $E$ 中有极限。

% SOURCE p=30; printed=15; object=ch00-proof-compactness-b-a-start; region=bottom
\subcref{theorem:ch00-compactness-characterization}{item:ch00-compact-bolzano-weierstrass}推出\subcref{theorem:ch00-compactness-characterization}{item:ch00-compact-complete-total}：我们证明，若后者的任一条件不成立，则前者也不成立。若 $E$ 不完备，则存在 $E$ 中的 Cauchy 序列 $\{x_n\}$，它在 $E$ 中没有极限。$\{x_n\}$ 的任何子序列都不可能在 $E$ 中收敛，否则整个序列就会收敛到同一极限。另一方面，若 $E$ 不全有界，取 $\varepsilon>0$，
% SOURCE p=31; printed=16; object=ch00-proof-compactness-b-a-end; region=top
使 $E$ 不能被有限个半径为 $\varepsilon$ 的球覆盖。按如下方式归纳选取 $x_n\in E$：先任取 $x_1\in E$；已选定 $x_1,\ldots,x_n$ 后，取 $x_{n+1}\in E\setminus\bigcup_1^nB(\varepsilon,x_j)$。于是对所有 $n,m$，$\rho(x_n,x_m)>\varepsilon$，所以 $\{x_n\}$ 没有收敛子序列。

% SOURCE p=31; printed=16; object=ch00-proof-compactness-ab-c; region=upper
\subcref{theorem:ch00-compactness-characterization}{item:ch00-compact-complete-total}和\subcref{theorem:ch00-compactness-characterization}{item:ch00-compact-bolzano-weierstrass}推出\subcref{theorem:ch00-compactness-characterization}{item:ch00-compact-heine-borel}：只需证明，若\subcref{theorem:ch00-compactness-characterization}{item:ch00-compact-bolzano-weierstrass}成立，且 $\{V_\alpha\}_{\alpha\in A}$ 是 $E$ 的开覆盖，则存在 $\varepsilon>0$，使每个与 $E$ 相交的半径为 $\varepsilon$ 的球都包含于某个 $V_\alpha$，因为由\subcref{theorem:ch00-compactness-characterization}{item:ch00-compact-complete-total}，$E$ 可被有限个这样的球覆盖。反设，对每个 $n\in\mathbb{N}$，存在半径为 $2^{-n}$ 的球 $B_n$，使 $B_n\cap E\ne\varnothing$，且 $B_n$ 不包含于任何 $V_\alpha$。选取 $x_n\in B_n\cap E$；取一个子序列后，可设 $\{x_n\}$ 收敛于某个 $x\in E$。对某个 $\alpha$ 有 $x\in V_\alpha$；因为 $V_\alpha$ 为开集，存在 $\varepsilon>0$ 使 $B(\varepsilon,x)\subseteq V_\alpha$。但当 $n$ 充分大，使 $\rho(x_n,x)<\varepsilon/3$ 且 $2^{-n}<\varepsilon/3$ 时，就有 $B_n\subseteq B(\varepsilon,x)\subseteq V_\alpha$，与对 $B_n$ 的假设矛盾。

% SOURCE p=31; printed=16; object=ch00-proof-compactness-c-b; region=middle
\subcref{theorem:ch00-compactness-characterization}{item:ch00-compact-heine-borel}推出\subcref{theorem:ch00-compactness-characterization}{item:ch00-compact-bolzano-weierstrass}：若 $\{x_n\}$ 为 $E$ 中没有收敛子序列的序列，则对每个 $x\in E$，存在以 $x$ 为中心的球 $B_x$，它只包含有限多个指标对应的 $x_n$（否则某个子序列将收敛于 $x$）。于是 $\{B_x\}_{x\in E}$ 是 $E$ 的开覆盖，却没有有限子覆盖。
\end{proof}

% SOURCE p=31; printed=16; object=ch00-compact-definition; region=middle

\begin{definition}[紧集]\label{definition:ch00-compact-definition}
具有\subcref{theorem:ch00-compactness-characterization}{item:ch00-compact-complete-total}至\subcref{theorem:ch00-compactness-characterization}{item:ch00-compact-heine-borel}所述性质的集合 $E$ 称为\textbf{紧集}。
\end{definition}

\begin{remark}
每个紧集都是闭集（由\cref{proposition:ch00-closed-complete}）且有界；逆命题一般不成立，但在 $\mathbb{R}^n$ 中成立。
\end{remark}


% SOURCE p=31; printed=16; object=ch00-euclidean-compactness; number=0.26; region=middle
\begin{proposition}\label{proposition:ch00-euclidean-compactness}
$\mathbb{R}^n$ 中每个闭且有界的子集都是紧集。
\end{proposition}

% SOURCE p=31; printed=16; object=ch00-proof-euclidean-compactness; region=middle-lower
\begin{proof}
由于 $\mathbb{R}^n$ 的闭子集完备，只需证明 $\mathbb{R}^n$ 的有界子集全有界。每个有界集都包含于某个立方体
\begin{align*}
Q=[-R,R]^n=\{x\in\mathbb{R}^n:\max(|x_1|,\ldots,|x_n|)\leqslant R\},
\end{align*}
所以只需证明 $Q$ 全有界。给定 $\varepsilon>0$，取整数 $k>R\sqrt{n}/\varepsilon$，把区间 $[-R,R]$ 等分为 $k$ 份，由此将 $Q$ 表为 $k^n$ 个全等子立方体的并。这些子立方体的边长为 $2R/k$，所以直径为 $\sqrt{n}(2R/k)<2\varepsilon$，因而每个子立方体都包含于以其中心为球心、半径为 $\varepsilon$ 的球中。
\end{proof}

% SOURCE p=31; printed=16; object=ch00-equivalent-metrics; region=lower

\begin{definition}[等价度量]\label{definition:ch00-equivalent-metrics}
集合 $X$ 上的两个度量 $\rho_1$ 和 $\rho_2$ 称为\textbf{等价的}，是指
\begin{align*}
C\rho_1\leqslant\rho_2\leqslant C'\rho_1\qquad\text{对某些 }C,C'>0.
\end{align*}
\end{definition}


\begin{remark}
容易验证，等价度量定义相同的开集、闭集和紧集，相同的收敛序列和 Cauchy 序列，以及相同的连续映射和一致连续映射。因此，关于度量空间的大多数结果只依赖于度量的等价类，而不依赖于具体选取的度量。
\end{remark}

\end{document}
